\documentclass[10pt,a4paper]{amsart}
\usepackage[margin=19mm]{geometry}
\usepackage{amsmath,amssymb,mathtools}
\usepackage[scr=rsfs,cal=euler]{mathalfa}
\usepackage{xfrac}
\usepackage[unicode,hidelinks,bookmarks=false]{hyperref}
\newcommand{\ie}{i.e.\ }
\newcommand{\cf}{cf.\ }
\newcommand{\resp}{respectively\ }
\newcommand{\Subsection}{\subsection}
\providecommand{\nobreakdash}{\nobreak\hbox{-}}
\setlength{\emergencystretch}{2em}
\multlinegap0pt
\usepackage{amssymb}
\usepackage[shortlabels]{enumitem}
\usepackage{mathtools}
\usepackage{bm}
\newtheorem{lemma}{Lemma}
\newcommand{\bb}[1]{\mathbb{#1}}
\title{Diagonal dimension and intermediate sub-\texorpdfstring{$\mathrm{C}^\ast$}{C*}-algebras}
\author{Grigoris Kopsacheilis}
\date{Source: arXiv:2607.24699v1 (CC BY 4.0)}
\begin{document}
\maketitle

\noindent\emph{Source and licence.} Diagonal dimension and intermediate sub-\texorpdfstring{$\mathrm{C}^\ast$}{C*}-algebras. Version arXiv:2607.24699v1 (CC BY 4.0), distributed by arXiv under the Creative Commons Attribution 4.0 International licence, http://creativecommons.org/licenses/by/4.0/, as recorded in the arXiv licence metadata for this exact version and shown on the abstract page https://arxiv.org/abs/2607.24699v1. Full licence notice: \texttt{licenses/arxiv-2607.24699-CC-BY-4.0.txt}.\par\smallskip

\noindent\textit{Editorial scope.} Counted unit: the article's lemma lem:orthogonal-range-support-0exp (source0.tex lines 213--215). The article's complete original proof of it is delivered with it (source0.tex lines 217--232, one \texttt{proof} environment of 16 lines). No mathematics is written, repaired or re-derived: the statement and the proof are the article's own lines, in the article's own notation, and no retained line is substituted or re-flowed.  No further source-local context is needed: every cross-reference of every retained line resolves inside the two delivered spans.  Nothing was omitted from any retained span, so the package carries no omission record.  No retained line contains a citation, so the wrapper carries no bibliography and no bibliography entry.  No retained line contains a figure environment, an image-inclusion command or drawing code, so no image is delivered and the package declares no asset dependency.  Editorial formatting only, in addition: an AMS \texttt{amsart} preamble that restates the article's own declarations and macros used by the retained lines, namely the environments \texttt{lemma} on separate counters, together with the standard packages those lines need.  The retained environments print in this document as Lemma 1 in order of appearance, whereas the article prints the same items with the numbering of its own full document.  The complete native source of the article as distributed in the arXiv e-print for this exact version is bundled unchanged as \texttt{sources/206138/source0.tex}.
\begin{lemma}\label{lem:orthogonal-range-support-0exp}
Let $(D\subset A)$ be an abelian sub-$\mathrm{C}^\ast$-algebra and let $\Phi\colon A\to D$ be a conditional expectation. If $v\in A$ is such that $vv^*, v^*v\in D$ and $vv^*\perp v^*v$, then $\Phi(v)=0$.
\end{lemma}

\begin{proof}
Note that $v=\lim_nv(v^*v)^\frac{1}{n}=\lim_n (vv^*)^\frac{1}{n}v$. For each $n\ge1$ let $\{p_{k,n}\}_{k\in\bb{N}}$ be a sequence of polynomials on $\bb{C}$ with zero constant term that converge uniformly over $\sigma(v^*v)\cup\{0\}=\sigma(vv^*)\cup\{0\}$ to the map $z\mapsto z^\frac{1}{n}$ (which exists by Weierstra\ss 's approximation theorem). Using that $D$ lies in the multiplicative domain of $\Phi$, we have that
\begin{align*}
\Phi(v)&=\lim_n\Phi(v(v^*v)^\frac{1}{n})\\
&=\lim_n\Phi(v)\cdot (v^*v)^\frac{1}{n}\\
&=\lim_n\lim_k\Phi(v)\cdot p_{k,n}(v^*v)
\end{align*}
so it suffices to show that $\Phi(v)v^*v=0$. But indeed,
\begin{align*}
\Phi(v)v^*v&=\lim_n\Phi((vv^*)^\frac{1}{n}v)v^*v\\
&=\lim_n(vv^*)^\frac{1}{n}\Phi(v)v^*v\\
&=\lim_n(vv^*)^\frac{1}{n}\cdot v^*v\cdot \Phi(v)\\
&=0
\end{align*}
where in the last step we used that $v^*v\perp vv^*$.
\end{proof}

\end{document}
